% Add "answers" option to print the answers
\documentclass[addpoints]{exam}

\usepackage[utf8]{inputenc}
\usepackage{calc}
\usepackage[dvipsnames]{xcolor}
\usepackage{ifthen}
\usepackage{tcolorbox}

% \makebox[\textwidth]{Name:\enspace\hrulefill}
%
% \vspace{5mm}
%\makebox[\textwidth]{Instructor’s name:\enspace\hrulefill}

\definecolor{MyLightGray}{gray}{0.60}

\newcommand{\truefalse}
{		\begin{checkboxes}
			 \choice True
			 \choice False. Explanation:
		\end{checkboxes}
		\vspace{10pt}
}


% Check if 'answers' is an option
\newcommand{\ifanswers}[2]{%
  \ifthenelse{\boolean{exam@answers}}{#1}{#2}%
}



\newcommand{\answerspace}[1]
{
	\ifprintanswers
	\else
		\setlength{\mylength}{10pt}%
		% \vspace{#1\mylength}
		% Make the grey box
		\par\nobreak\noindent
		\makebox[\linewidth][l]{%
			% set inner padding to zero, so only the frame shows
			{\setlength{\fboxsep}{0pt}%
			% draw a color frame, white (or transparent) inside
			\fcolorbox{MyLightGray}{white}{%
			% the inner box has zero width content but fixed height
			\parbox[t][#1\mylength][c]{%
				% subtract twice the frame‐thickness from the full width
				\dimexpr\linewidth-2\fboxrule\relax
			}{\strut}%
			}%
			}%
		}%
		\par
		\vspace{3\mylength}

  	\fi
}
\newlength{\mylength}

\newcommand{\code}[1]{{\texttt{#1}}}


\usepackage{color}
\usepackage{lastpage}
\usepackage{verbatim}
\usepackage{fancyvrb}
\usepackage{graphicx}
\usepackage{float}
\usepackage{minted}
\usepackage{multicol}
\usepackage{amsmath}

% For those sweet, sweet acmart fonts.
\usepackage[T1]{fontenc}
\usepackage[tt=false, type1=true]{libertine}
\usepackage[varqu]{zi4}                  % Inconsolata for \ttfamily

\RecustomVerbatimCommand{\VerbatimInput}{VerbatimInput}{}

\newcommand{\Code}[1]{\mintinline{python}{#1}}

\renewcommand{\solutiontitle}{\noindent\textbf{Solution sketch: }}
\definecolor{SolutionColor}{rgb}{0.8,0.9,1}

\cfoot{Page \thepage\ of \pageref{LastPage}}
\rhead{Initials: \rule{20mm}{0.4pt}}

%%%%%%%%%%%%%%%%%%%%%%%%%%%%%%%%%%%%%%%%
\begin{document}
\fvset{numbers=left,numbersep=3pt}
\thispagestyle{empty}

\begin{center}
  {\LARGE CS 408 Fall 2026 -- Concolic Execution Quiz \\}
    \vspace*{5pt}
\fbox{\fbox{\parbox{5.5in}{\centering
	\begin{itemize}
        \item Try to come up with a solution on your own first, but
          after that talk to your friends and neighbors.

        \item Full points will be rewarded for \textbf{any} (reasonable)
          attempt.
        \end{itemize}
	}}}
\end{center}

\vspace{5mm}
\makebox[\textwidth]{\textbf{Name (please print)}: \rule{100mm}{0.8pt} }

\begin{questions}

\question Here is the \Code{factorial} function from Tuesday's code demo:

\begin{minted}{python}
def factorial(n: int) -> int:
    v = 1
    while n != 0:
        v = v * n
        n = n - 1
    return v
\end{minted}

  Running it concolically on \Code{n = 2} produces the following path
  constraints, where \Code{a} is a symbolic value:
  \begin{center}
    $a \neq 0 \;\wedge\; a - 1 \neq 0 \;\wedge\; \neg(a - 1 - 1 \neq 0)$
  \end{center}

\begin{parts}

\part How does the size of the path condition change as \Code{n} gets
  larger?  Roughly how many constraints would \Code{factorial(10)}
  produce? %
  \answerspace{14} %
  \begin{solution} %
    Each trip through the loop tests \Code{n != 0} once, and that test
    contributes one constraint. Running on \Code{n = 10} executes the
    test eleven times --- ten that succeed and one that fails --- so the
    path condition has \textbf{11} constraints.

    Every distinct trip count is a distinct path. A loop whose bound is
    an input has as many paths as the input has possible values, which
    is why loops are the hard case for both symbolic and concolic
    execution.
  \end{solution}

%\newpage

\part To reach a new path, the CE engine takes a \textbf{prefix} of
  the path condition and negates its last constraint. What is one
  possible input that an SMT solver could return if we used a prefix
  of the above path condition containing just the first constraint? %
  \answerspace{9} %
  \begin{solution} %
    The only solution is \Code{n = 0}. The loop body never runs,
    \Code{factorial(0)} returns 1, and the resulting path condition is
    the single constraint $\neg(N \neq 0)$.

    That is the other side of the very first branch --- one new path for
    one solver query, which is exactly what concolic execution is for.
  \end{solution}

\newpage

\part Suppose that the CE engine used the first two constraints as the
  prefix instead. What are 2-3 possible inputs that the solver might
  generate using this this prefix? Could any of these trigger a
  (potential) failure in \Code{factorial}? %
  \answerspace{18} %
  \begin{solution} %
    The input is \textbf{correct}: $-1 \neq 0$ holds, and
    $-1 - 1 = -2 \neq 0$ holds. The solver did its job.

    But \Code{factorial(-1)} \textbf{never terminates}. The loop
    decrements: $-1, -2, -3, \ldots$ and never reaches 0, so the guard
    \Code{n != 0} is true forever.

    (This really happens --- it is what hung the demo in class.)
  \end{solution}

% \newpage

% \part The solver produced an input that satisfies every constraint it
%   was given, and the program still misbehaved. What did the path
%   condition fail to capture, and what does a real tool do about it? %
%   \answerspace{20} %
%   \begin{solution} %
%     A path condition records only the branches that were \emph{actually
%     taken on the observed run}. On \Code{n = 2} the loop always counted
%     down toward zero, so nothing in the program ever tested whether
%     \Code{n} was positive --- and so no constraint says so. The prefix
%     ``the loop runs at least twice'' is satisfied by $n = -1$ even
%     though the real program can never leave the loop from there.

%     This is \textbf{divergence}: the input solves for one path, and the
%     execution goes somewhere else entirely.

%     Real tools cope by bounding execution --- a timeout, or a cap on
%     loop iterations --- and discarding inputs that exceed it. Notice
%     what the bound costs: non-termination is not a crash and not a
%     failed assertion, so \emph{no oracle we have seen this semester
%     detects it}. We are reduced to guessing how long is too long.
%   \end{solution}

\end{parts}

\end{questions}

\end{document}
